\section{Why It Runs Once Per Session}
\label{sec:ch03-why-it-runs-once-per-session}

\index{execution!field by field|(}%
The engine resolved \code{sessions} first and got four objects back. Then it
had a list of four and a field to resolve on each of them, and it did the only
thing it can do without knowing anything about your resolver: it called the
resolver four times.

That is not a defect and there is no setting that turns it off. A resolver is a
function from one parent object to one value, and the four calls are the
literal meaning of asking for a field on a list of four things. The engine
cannot merge them because it cannot see inside them. As far as it knows,
\code{GetSpeakerAsync} reads a file, calls a payment provider or returns a
constant, and nothing in its signature says which.

\index{DataLoader!what it changes about timing}%
What can be changed is when the four calls turn into work. The engine knows
something your resolver does not: it knows that all four are happening at the
same moment, on the same list, in the same request. If the resolver could hand
back a promise instead of a query, and the four promises could be collected
before any of them was answered, the four lookups become one lookup for four
keys. That is the entire idea, and the thing that implements it is called a
DataLoader.

Figure~\ref{fig:ch03-one-batch} is the two arrangements side by side. The
difference is not that lane B does less work. It is that lane B does its
database work after the engine has finished asking, and lane A does its work
during.

\begin{figure}
  \centering
  % Four conference sessions, each with a speaker, resolved two ways: once a
% resolver at a time, once behind a batching data loader.
%
% Same idiom as figures/tikz/ch01-one-field.tex and ch02-two-requests.tex:
% each lane is a timeline read left to right, ticks mark stages, and a brace
% under the stages that matter carries a one-line label. Lane A and lane B
% share the same sessions resolver and the same first statement to SQLite,
% so those sit at the same x in both lanes; only what happens after that
% diverges. Lane A repeats the speaker resolver, and a database statement,
% once per session, drawn as four separate ticks rather than an ellipsis, so
% the interleaving with SQLite reads on the page. Lane B calls the speaker
% resolver the same four times but only records a key each time; a dashed
% boundary marks where the engine stops asking, then one statement carries
% the whole batch. Four sessions have three distinct speakers, so the batch
% carries three keys, not four; the duplicate collapses before it reaches
% SQLite.
%
% Bare tikzpicture (house style): the figure environment, caption and label
% live at the call site in the section file.
\begin{tikzpicture}[
  x=1mm, y=1mm,
  every node/.append style={font=\sffamily\scriptsize},
  axis/.style={draw, line width=0.4pt, -{Straight Barb[length=1.4mm, width=1.6mm]}},
  tick/.style={draw, line width=0.4pt},
  event/.style={anchor=south, align=center, text width=22mm, inner sep=1pt},
  span/.style={draw, line width=0.4pt},
  spanlabel/.style={anchor=north, align=center, inner sep=2pt},
  lane/.style={anchor=west, font=\sffamily\scriptsize\itshape},
  boundary/.style={draw, line width=0.4pt, dashed},
]

% ------------------------------------------------------------------- lane A
\node[lane] at (0,13) {A. one resolver call per session, each reaching SQLite};

\draw[axis] (0,0) -- (150,0);
\foreach \x in {15,40,60,80,100,135}{\draw[tick] (\x,-1.5) -- (\x,1.5);}

\node[event, text width=24mm] at (15,3)  {sessions resolver:\\stmt 1, 4 sessions};
\node[event, text width=16mm] at (40,3)  {session 1\\stmt 2};
\node[event, text width=16mm] at (60,3)  {session 2\\stmt 3};
\node[event, text width=16mm] at (80,3)  {session 3\\stmt 4};
\node[event, text width=16mm] at (100,3) {session 4\\stmt 5};
\node[event, text width=20mm] at (135,3) {response\\assembled};

\draw[span] (32,-4) -- (32,-7) -- (108,-7) -- (108,-4);
\node[spanlabel] at (70,-8) {speaker resolver, once per session:\\statements 2 to 5};

% ------------------------------------------------------------------- lane B
\begin{scope}[shift={(0,-46)}]
  \node[lane] at (0,13) {B. the same four calls, one batched statement};

  \draw[axis] (0,0) -- (150,0);
  \foreach \x in {15,108,135}{\draw[tick] (\x,-1.5) -- (\x,1.5);}
  \foreach \x in {40,46,52,58}{\draw[tick] (\x,-1.5) -- (\x,1.5);}
  \draw[boundary] (84,-3) -- (84,6);

  \node[event, text width=24mm] at (15,3)  {sessions resolver:\\stmt 1, 4 sessions};
  \node[event, text width=32mm] at (49,3)  {speaker resolver x4:\\records key, no query};
  \node[event, text width=22mm] at (108,3) {one statement:\\3 distinct keys};
  \node[event, text width=27mm] at (134,3) {promises completed:\\response assembled};

  \draw[span] (36,-4) -- (36,-7) -- (62,-7) -- (62,-4);
  \node[spanlabel] at (49,-8) {no statement issued};

  \node[spanlabel] at (84,-8) {batch boundary:\\engine stops asking};

  \draw[span] (100,-4) -- (100,-7) -- (116,-7) -- (116,-4);
  \node[spanlabel] at (108,-8) {1 statement,\\3 distinct keys};
\end{scope}

\end{tikzpicture}

  \caption{The same four sessions, resolved two ways. In lane A each call to
  the speaker resolver reaches SQLite on its own, so the statements interleave
  with the calls and there are as many as there are sessions. In lane B the
  resolver returns a promise and records a key; nothing reaches SQLite until
  the engine has stopped asking, at which point three distinct keys go out in
  one statement. The batch boundary, not the amount of data, is what the
  DataLoader moves.}
  \label{fig:ch03-one-batch}
\end{figure}

\index{DataLoader!deduplicates keys}%
Lane B also shows the second thing a DataLoader does, which is easy to miss
because it is invisible in the code. Four sessions produced three keys, because
Ada Fischer gives two of them and a set does not hold a duplicate. The naive
resolver could not have known that, since each of its calls saw one session and
had no way to ask what the other three wanted.

\index{N+1@N+1!across a seam}%
Hold on to the shape of lane A rather than the fix. In this chapter the four
repeated statements go to a database in the same process, and the worst that
happens is that you notice. In chapter~\ref{ch:router-n-plus-1} the same
picture has a router in the place of the engine and an HTTP request in the
place of each statement, the parent list arrives from one service and the
field belongs to another, and the four calls are four round trips over a
network you do not control. The mechanism is identical, and the unit of
repetition stops being a local statement and becomes a request to another
team's service.
\index{execution!field by field|)}%
